\chapter{Diseño del sistema}
\label{cap:diseno-sistema}

\section{Visión general metodológica}
Este capítulo describe el diseño metodológico empleado para construir y evaluar el defensor de ciberseguridad basado en aprendizaje por refuerzo (RL) propuesto. El objetivo no es presentar un sistema de prevención de intrusiones (IPS) listo para producción, sino definir un marco experimental controlado, reproducible y offline para decisiones de seguridad binarias sobre registros de flujos de red. Cada flujo se trata como un único punto de decisión: el defensor observa una representación numérica del flujo y decide si permitirlo (\texttt{PERMIT}) como tráfico benigno o bloquearlo (\texttt{BLOCK}) como tráfico malicioso.

La metodología es deliberadamente conservadora. Los benchmarks públicos de detección de intrusiones pueden producir resultados inflados cuando el preprocesamiento, la construcción de las particiones, los controles de fuga de información y las comparaciones con líneas base son débiles \parencite{Arp2020DosDontsMLSecurity,Engelen2021CICIDSIssues,Lanvin2023Errors}. Por ello, el pipeline se ha diseñado en torno a contratos de datos explícitos, artefactos de ejecución reproducibles, normalización ajustada únicamente sobre datos de entrenamiento, y una validación progresivamente más estricta. La misma representación canónica se utiliza tanto para el benchmark interno de CICIDS2017 como para la inferencia offline sobre tráfico de laboratorio en la Fase 2, de modo que la pregunta experimental no queda ligada a una convención de nomenclatura de CSV concreta ni a un extractor de características particular.

A grandes rasgos, la metodología comprende seis etapas encadenadas:
\begin{enumerate}
    \item ingerir los datos CSV de CICIDS2017 a nivel de flujo y mapear las etiquetas a una tarea binaria de defensa;
    \item eliminar los identificadores susceptibles de provocar fuga de información y normalizar los valores numéricos malformados;
    \item mapear las columnas específicas de cada extractor a un esquema canónico de características fijo de 76 características;
    \item añadir una máscara de presencia/ausencia de 76 dimensiones, produciendo observaciones de 152 dimensiones;
    \item entrenar y evaluar un defensor QRDQN compatible con Gymnasium bajo protocolos de partición controlados;
    \item comprobar la robustez mediante comprobaciones de validación, líneas base supervisadas e inferencia offline en la Fase 2.
\end{enumerate}

Este diseño mantiene la tesis dentro de un marco experimental offline. Las salidas \texttt{PERMIT} y \texttt{BLOCK} son predicciones del modelo sobre registros de flujo almacenados. No son órdenes de cortafuegos y no modifican ninguna ruta de red activa.

\section{Fuente de datos y representación de la entrada}
La tesis utiliza una representación tabular a nivel de flujo. Esta representación es un compromiso deliberado entre relevancia operativa, manejabilidad y privacidad. Captura el comportamiento estadístico a lo largo de secuencias de paquetes evitando la inspección de la carga útil.

\subsection{CSVs de flujos de CICIDS2017}
La fuente de datos principal es la publicación de CSVs de flujos de CICIDS2017. Cada fila de un CSV corresponde a un flujo de red bidireccional resumido mediante estadísticas numéricas como duración, conteos de paquetes, totales de bytes, tiempos de llegada entre paquetes, estadísticas de longitud de paquete y conteos de flags TCP. Las representaciones basadas en flujos se utilizan ampliamente en la investigación sobre NIDS porque comprimen el comportamiento de la red en filas numéricas de anchura fija que pueden procesarse mediante métodos estándar de aprendizaje automático \parencite{Sperotto2010FlowIDS,Umer2017FlowIDS}.

El pipeline de carga local espera los ficheros CSV oficiales de CICIDS2017 y admite presupuestos de filas tanto completos como limitados. Los presupuestos limitados permiten pruebas rápidas de humo e iteración ágil; los presupuestos completos son los utilizados en las ejecuciones finales del benchmark. El cargador lee los CSVs grandes en fragmentos, estandariza los nombres de columna eliminando espacios en blanco inconsistentes, identifica la columna de etiqueta de forma robusta y aplica después la limpieza de características antes de cualquier escalado dependiente de la partición.

\subsection{Mapeo de etiquetas binarias}
CICIDS2017 incluye varias etiquetas de ataque. La tesis las colapsa en una decisión binaria del defensor:
\begin{itemize}
    \item etiqueta \texttt{0}: \texttt{BENIGN}, mapeada a la acción deseada \texttt{PERMIT};
    \item etiqueta \texttt{1}: cualquier clase de ataque, mapeada a la acción deseada \texttt{BLOCK}.
\end{itemize}

Este mapeo binario se corresponde con el modelo de decisión de primera línea estudiado aquí. Además, hace que el espacio de acciones de RL sea sencillo y auditable. La limitación es que la información sobre la familia de ataque no es el objetivo de aprendizaje principal en la ruta de RL actual. Cualquier análisis de errores por familia debe por tanto estar respaldado por etiquetas de familia preservadas o artefactos independientes; no puede inferirse a partir de las etiquetas binarias por sí solas.

\subsection{Representación tabular a nivel de flujo}
El sustrato de aprendizaje es una tabla a nivel de flujo en lugar de paquetes en bruto o bytes de carga útil. Esto reduce la dimensionalidad y hace factible el entrenamiento offline a gran escala. También evita algunas complicaciones legales y técnicas de la inspección profunda de paquetes, especialmente cuando las cargas útiles están cifradas \parencite{Sperotto2010FlowIDS,Umer2017FlowIDS}. Sin embargo, la abstracción de flujo pierde el contenido, el orden preciso de los paquetes y parte del contexto de largo alcance entre múltiples flujos. La metodología estudia por tanto un defensor basado en flujos, no un sistema completo de análisis forense de red.

\section{Limpieza de datos y preprocesamiento}
El preprocesamiento se trata como parte del método experimental, no como un detalle de implementación. En los benchmarks de NIDS, pequeñas decisiones de preprocesamiento pueden cambiar el significado de la tarea al filtrar etiquetas, suprimir flujos malformados o hacer que los datos de prueba influyan en el entrenamiento.

\subsection{Normalización de columnas y coerción numérica}
Las columnas CSV en bruto pueden contener espacios iniciales, capitalización inconsistente y tipos de datos mixtos. El cargador primero normaliza los nombres de columna y después convierte las columnas de características a valores numéricos. La matriz de características final se representa como \texttt{float32}, mientras que las etiquetas se representan como \texttt{int64}. Este contrato de tipos mantiene estable la interfaz del entorno, el escalador y el modelo en las etapas posteriores.

No se permite que las filas introduzcan cadenas de texto arbitrarias en el modelo. Las columnas de características no numéricas que sobreviven a la eliminación de identificadores quedan excluidas. Esto evita la introducción accidental de cadenas categóricas cuya codificación sería no controlada o específica de un conjunto de datos.

\subsection{Valores inválidos, valores ausentes e imputación}
La extracción de flujos de red puede producir valores numéricos inválidos. Las características de tasa, por ejemplo, pueden volverse infinitas o indefinidas cuando se calculan sobre una duración nula o casi nula. El pipeline sustituye los infinitos y los valores numéricos ausentes por cero antes del mapeo canónico. Esta elección es pragmática: muchos contadores de flujo, tasas y estadísticas agregadas pueden utilizar de forma segura el cero como marcador de posición neutro cuando se acompaña de una señal explícita de presencia/ausencia.

La imputación por cero sola sería insegura porque el modelo no podría distinguir un cero medido real de un marcador sintético. El esquema canónico empareja por tanto la imputación con una máscara de presencia/ausencia. La máscara hace visible la disponibilidad de características tanto para el modelo como para los diagnósticos posteriores.

\subsection{Exclusión de campos susceptibles de fuga}
El pipeline de preprocesamiento aplica una política anti-fuga fundamentada en los riesgos de evaluación conocidos en seguridad con ML \parencite{Arp2020DosDontsMLSecurity,Kaufman2011LeakageDataMining}. Se excluyen las direcciones IP de origen y destino, los Flow IDs, las marcas temporales absolutas y los identificadores de endpoints. Los campos de puerto directo también se excluyen porque pueden actuar como indicadores del escenario en conjuntos de datos controlados por scripts. Por ejemplo, un día de ataque puede concentrar tráfico malicioso en servicios concretos, lo que permitiría al modelo memorizar un script de captura en lugar de aprender estadísticas de comportamiento de flujo.

Esta exclusión se aplica antes del mapeo canónico. El modelo recibe por tanto únicamente características estadísticas de flujo y los valores de máscara correspondientes. Si en el futuro se introducen nuevos conjuntos de datos o extractores, sus adaptadores deben preservar la misma política: sin identificadores, sin campos de tiempo absoluto y sin puertos utilizados directamente como atajos de etiqueta.

\section{Esquema canónico de características}
El esquema canónico es el contrato de datos central de la tesis. Separa la interfaz del modelo de los nombres de columna, el orden y la disponibilidad de características de un extractor concreto.

\subsection{Características canónicas de flujo}
El esquema define 76 características numéricas de flujo ordenadas. Incluyen duración, conteos de paquetes hacia adelante y hacia atrás, totales de bytes, estadísticas de longitud de paquete, estadísticas de tiempo entre llegadas de flujo y por dirección, tasas de paquetes y bytes, conteos de flags TCP, longitudes de cabecera, tamaños de ventana, estadísticas de subflujos y características de temporización activa/inactiva. Estas características se han elegido porque son estadísticas de flujo y no identificadores, y porque son ampliamente compatibles con la extracción al estilo CICFlowMeter.

Para CICIDS2017, un adaptador mapea las columnas originales del CSV a estos nombres canónicos. Para la Fase 2, el script de inferencia robusta mapea las columnas del extractor de flujos de laboratorio a los mismos nombres canónicos. El modelo se entrena y evalúa por tanto frente a un contrato de características fijo. Este diseño sigue la motivación metodológica más amplia de la estandarización de características en NIDS: la comparación entre dominios resulta más auditable cuando las definiciones de características son explícitas y estables \parencite{Sarhan2021}.

La Tabla~\ref{tab:canonical-groups} enumera los grupos lógicos dentro del esquema de 76 características y el número de características aportadas por cada grupo. La agrupación es informativa; el modelo recibe los 76 valores como un vector plano sin señales de agrupación estructural.

\begin{table}[htbp]
\centering
\begin{tabular}{lcp{7cm}}
\toprule
Grupo de características & Cantidad & Estadísticas representativas \\
\midrule
Duración del flujo & 1 & Duración total del flujo bidireccional \\
Estadísticas de paquetes hacia adelante & 12 & Conteo de paquetes, bytes totales, media/desv.\ típica/mínimo/máximo de longitudes de paquete \\
Estadísticas de paquetes hacia atrás & 12 & Las mismas estadísticas para la dirección inversa \\
Agregados bidireccionales & 8 & Conteos combinados, bytes totales, promedios de tasa de ráfaga \\
Estadísticas de tiempo entre llegadas & 10 & Media, desv.\ típica, mínimo y máximo del tiempo entre llegadas del flujo, hacia adelante y hacia atrás \\
Conteos de flags TCP & 8 & Conteos de FIN, SYN, RST, PSH, ACK, URG, CWE, ECE \\
Agregados de longitud de cabecera & 4 & Totales de bytes de cabecera hacia adelante y hacia atrás \\
Estadísticas de tamaño de ventana & 3 & Tamaño de ventana inicial y medio por dirección \\
Características de tasa de flujo & 6 & Bytes/s, paquetes/s por dirección y combinados \\
Estadísticas de subflujos & 6 & Promedios de paquetes y bytes de subflujo por dirección \\
Temporización activa/inactiva & 6 & Media, desv.\ típica y máximo de las duraciones de los periodos activos e inactivos \\
\bottomrule
\end{tabular}
\caption[Agrupación lógica de las 76 características canónicas]{Agrupación lógica de las 76 características canónicas de flujo. Todos los grupos se fusionan en un único vector plano antes de la entrada al modelo.}
\label{tab:canonical-groups}
\end{table}

\subsection{Máscara de presencia/ausencia}
Diferentes extractores pueden no proporcionar todas las características canónicas, e incluso una columna fuente presente puede contener valores inválidos. El pipeline produce por tanto una máscara de presencia/ausencia de 76 dimensiones. Para cada característica canónica $x_i$, el valor de máscara correspondiente $m_i$ es:
\begin{itemize}
    \item $m_i = 1$ cuando la característica fuente está presente y es válida;
    \item $m_i = 0$ cuando la característica está ausente o ha tenido que ser imputada.
\end{itemize}

La máscara desempeña dos funciones metodológicas. En primer lugar, evita la imputación silenciosa: el modelo puede condicionar su decisión a si un valor fue observado o rellenado. En segundo lugar, facilita los diagnósticos de la Fase 2, porque un extractor de laboratorio que carece de muchas características al estilo CICIDS2017 puede detectarse mediante un patrón de máscara diferente en lugar de quedar oculto dentro de un vector numérico denso.

Conviene precisar la semántica exacta de la máscara en la implementación actual. La máscara se calcula \emph{después} de la limpieza que sustituye los valores no finitos por cero (\texttt{src/load\_cicids2017.py}: $\pm\infty \to$ NaN $\to 0$), de modo que $m_i = 1$ siempre que la característica canónica tiene una columna fuente mapeada. Como el mapeo CICIDS2017 $\to$ canónico cubre las 76 características, sobre los datos nativos de CICIDS2017 la máscara es \textbf{constante e igual a 1} en todas las filas: codifica la \emph{presencia de la columna fuente}, no la ausencia de valores fila a fila. La máscara solo se vuelve informativa en la inferencia entre dominios o en la Fase 2 de laboratorio, donde algunas columnas fuente no están disponibles y su valor de máscara permanece en 0; un mapeo histórico como NSL-KDD, por ejemplo, solo cubre 3 de las 76 características canónicas. Así, las 152 dimensiones de la observación se componen de 76 valores de características más 76 indicadores de presencia.

\subsection{Vector de observación final}
La observación final es:
\[
    o = [x_1, x_2, \dots, x_{76}, m_1, m_2, \dots, m_{76}]
\]
Esto produce un vector \texttt{float32} de 152 dimensiones. La primera mitad contiene los valores de las características, y la segunda mitad contiene la máscara. El tamaño de la observación es por tanto invariante entre CICIDS2017, las particiones de validación y la inferencia en la Fase 2. Esta invariancia es importante porque la red de política del QRDQN espera una dimensión de entrada fija.

\section{Formulación del aprendizaje por refuerzo}
El defensor se implementa como un entorno compatible con Gymnasium \parencite{GymnasiumDocs}. Esto permite que el código estándar de RL basado en valores interactúe con el conjunto de datos de flujos a través de la interfaz habitual \texttt{reset()} y \texttt{step()}.

\subsection{Espacio de observación}
El espacio de observación es un espacio vectorial continuo de dimensión 152. Cada observación representa un vector de flujo canónico escalado junto con su máscara de presencia/ausencia. Durante el entrenamiento, las observaciones se extraen de la partición de entrenamiento. Durante la evaluación determinista, el entorno itera a través de la partición de prueba sin barajado, lo que permite comparar las predicciones directamente con las etiquetas verdaderas.

\subsection{Espacio de acciones}
El espacio de acciones es discreto con dos acciones:
\begin{itemize}
    \item \texttt{Acción 0 (PERMIT)}: clasificar el flujo como benigno y permitirlo en el modelo de decisión experimental;
    \item \texttt{Acción 1 (BLOCK)}: clasificar el flujo como malicioso y bloquearlo o generar una alerta en el modelo de decisión experimental.
\end{itemize}

El valor de la acción está alineado intencionalmente con el valor de la etiqueta binaria: benigno es \texttt{0}, ataque es \texttt{1}. Esto simplifica la evaluación directa porque una acción predicha puede compararse con la etiqueta objetivo sin necesidad de una capa de decodificación adicional.

\subsection{Función de recompensa}
La función de recompensa codifica costes de seguridad asimétricos \parencite{Lee2002CostSensitiveIDS,Cardenas2006IDSFramework}. Un falso negativo significa que se permite un ataque. Un falso positivo significa que se bloquea tráfico benigno. En muchos escenarios defensivos, los ataques no detectados son más graves que los bloqueos innecesarios, aunque la relación exacta depende del contexto operativo.

La implementación actual utiliza la siguiente configuración de recompensa por defecto:
\[
\text{TP}=1.5,\quad \text{FP}=-2.0,\quad \text{FN}=-5.0,\quad \text{TN/omisión}=0.0.
\]
La Tabla~\ref{tab:reward-matrix} organiza estos valores como una matriz de decisión para los cuatro resultados posibles.

\begin{table}[htbp]
\centering
\begin{tabular}{lcc}
\toprule
Decisión del agente & Etiqueta real: Ataque & Etiqueta real: Benigno \\
\midrule
\texttt{BLOCK} (Acción 1) & $+1.5$ (Verdadero Positivo) & $-2.0$ (Falso Positivo) \\
\texttt{PERMIT} (Acción 0) & $-5.0$ (Falso Negativo) & $\phantom{-}0.0$ (Verdadero Negativo) \\
\bottomrule
\end{tabular}
\caption[Matriz de recompensa del entorno del defensor binario]{Matriz de recompensa por defecto para el entorno del defensor binario. Los valores son supuestos experimentales registrados con cada ejecución.}
\label{tab:reward-matrix}
\end{table}

La penalización por falso negativo domina la señal de recompensa. Un agente que emite \texttt{PERMIT} en cada observación acumula un retorno muy negativo en tráfico con alta proporción de ataques, mientras que un agente que emite \texttt{BLOCK} en cada observación acumula penalizaciones por falsos positivos sin obtener crédito por detección de ataques. La estructura de recompensa está por tanto diseñada para alejar al agente de políticas constantes degeneradas y conducirlo hacia un comportamiento discriminativo. La recompensa cero para el verdadero negativo es deliberada: recompensar explícitamente cada paso benigno correcto sesgaría al agente hacia maximizar la tasa de permiso en lugar de la calidad de detección. Estos valores son supuestos experimentales, no costes empresariales medidos. Se registran con cada ejecución para que las ejecuciones históricas puedan distinguirse de los valores por defecto actuales.

\subsection{Estructura del episodio}
Un episodio es una secuencia de decisiones sobre flujos. Al reiniciar, el entorno inicializa un índice sobre la partición seleccionada. Cuando el barajado está activado, el orden de entrenamiento se aleatoriza. En cada paso, el agente recibe una observación, selecciona \texttt{PERMIT} o \texttt{BLOCK}, recibe una recompensa calculada a partir de la etiqueta real y la acción, y avanza al flujo siguiente. Un episodio termina cuando se agota el conjunto de datos o se alcanza un número máximo de pasos configurado.

Este diseño proporciona a los algoritmos estándar de RL un flujo de transiciones manteniendo el entrenamiento offline y seguro. Ninguna acción afecta al conjunto de datos fuente, al flujo siguiente ni al entorno de red.

\subsection{Estado del entorno y protocolo de transición}
En cada llamada a \texttt{reset()}, el entorno selecciona la partición asignada al modo actual (entrenamiento o evaluación), baraja opcionalmente el índice si el modo de barajado está activo, reinicia el contador de pasos interno y devuelve la primera observación junto con un diccionario de metadatos. En cada llamada a \texttt{step()}, el entorno registra la acción del agente, recupera la etiqueta de verdad del índice de flujo actual, calcula la recompensa escalar a partir de la matriz de recompensa, avanza el índice y devuelve la siguiente observación, la recompensa, un indicador de terminación, un indicador de truncación y el diccionario de metadatos.

La condición de terminación es el agotamiento del conjunto de datos o un número máximo de pasos configurable. En el modo de entrenamiento, el orden barajado garantiza que las transiciones del minibatch muestreen una distribución amplia de tipos de flujo durante un único episodio. En el modo de evaluación, el entorno es determinista: la misma partición recorrida en el mismo orden produce la misma secuencia de observaciones, recompensas y etiquetas. Este determinismo es esencial para la evaluación reproducible porque implica que las diferencias en métricas entre dos modelos guardados reflejan únicamente el comportamiento del modelo y no el ruido de muestreo en el bucle de evaluación.

La separación entre los modos de entrenamiento y evaluación se aplica a nivel del entorno y no se deja al llamante. Pasar el indicador de modo correcto en el momento de la construcción establece tanto el conjunto de índices como el comportamiento de barajado, lo que evita un error frecuente en el que la evaluación hereda accidentalmente el orden barajado de un entorno de entrenamiento.

\subsection{Limitaciones de la formulación conjunto-de-datos-como-entorno}
Dado que el entorno es un conjunto de datos estático, las acciones del agente no afectan a los estados futuros. Esto no es un problema completo de defensa cibernética autónoma en el que bloquear un flujo podría cambiar el comportamiento posterior del atacante, el estado del host, el enrutamiento o el contexto de alertas. La formulación se asemeja más a un problema de decisión contextual sensible al coste que a un proceso de decisión de Markov rico \parencite{SuttonBarto2018RL,Levine2020OfflineRL}.

Esta limitación es explícita en la metodología. El RL se utiliza aquí para estudiar el aprendizaje de decisiones basadas en valor y sensibles al coste sobre registros de flujo, y para establecer un pipeline offline seguro. Las afirmaciones sobre respuesta adversarial en múltiples pasos, adaptación en línea o defensa cibernética activa quedan fuera del alcance implementado.

\section{Agente QRDQN}
El algoritmo de RL principal es Quantile Regression Deep Q-Network (QRDQN) \parencite{Dabney2018QRDQN}. QRDQN pertenece a la familia del RL distribucional, que modela una distribución sobre los retornos en lugar de únicamente un retorno esperado \parencite{Bellemare2017DistributionalRL}.

\subsection{Papel algorítmico}
El DQN estándar aproxima los valores-acción con una red neuronal y selecciona acciones según el retorno esperado estimado \parencite{Mnih2015DQN}. QRDQN extiende esta idea aprendiendo los cuantiles de la distribución del retorno. Esto es relevante para la tesis porque la estructura de recompensa es asimétrica: los falsos negativos conllevan una penalización mucho mayor que la recompensa que generan los permisos benignos correctos. Una visión distribucional puede representar más información sobre la variabilidad del retorno que una única estimación de la media.

La metodología no asume que QRDQN sea inherentemente superior a los clasificadores supervisados para datos de flujo tabulares. En cambio, evalúa QRDQN como el representante principal de RL bajo el mismo preprocesamiento, las mismas particiones y las mismas métricas utilizadas por las líneas base.

\subsection{Función de pérdida de regresión cuantílica}
QRDQN representa la distribución del retorno de acción mediante $N$ estimaciones de cuantiles aprendidas $\{\theta_i(s, a)\}_{i=1}^{N}$ en puntos medios de cuantil fijos $\hat{\tau}_i = \frac{2i-1}{2N}$ \parencite{Dabney2018QRDQN}. Durante la selección de acción, estos cuantiles se promedian para producir una estimación del retorno esperado, recuperando la regla codiciosa estándar del DQN bajo expectativa.

Para una única transición $(s, a, r, s')$, sea $a^* = \arg\max_a \sum_j \theta_j^-(s', a)$ la acción codiciosa bajo la red objetivo, y sea
\[
    \delta_{ij} = r + \gamma\,\theta_j^-(s', a^*) - \theta_i(s, a)
\]
el error de diferencia temporal entre el cuantil fuente $i$ y el cuantil objetivo $j$ calculado mediante bootstrapping. QRDQN minimiza la función de pérdida de regresión cuantílica de Huber sumada sobre todos los $N^2$ pares de cuantiles:
\[
    \mathcal{L} = \frac{1}{N}\sum_{i=1}^{N}\sum_{j=1}^{N}
        \bigl|\hat{\tau}_i - \mathbf{1}[\delta_{ij} < 0]\bigr|\cdot L_\kappa(\delta_{ij}),
\]
donde $L_\kappa$ es la pérdida de Huber elemento a elemento con umbral $\kappa$:
\[
    L_\kappa(\delta) = \begin{cases} \tfrac{1}{2}\delta^2 & \text{si } |\delta| \leq \kappa, \\ \kappa\bigl(|\delta| - \tfrac{1}{2}\kappa\bigr) & \text{en caso contrario.}\end{cases}
\]
A diferencia de C51, que representa los retornos sobre un soporte categórico fijo, QRDQN no impone ninguna restricción sobre el soporte de la distribución del retorno y se adapta automáticamente a la escala de las recompensas \parencite{Bellemare2017DistributionalRL}. En el entorno de recompensa asimétrica de esta tesis, la distribución del retorno puede desarrollar colas negativas pronunciadas cuando los eventos de falso negativo son frecuentes, lo que hace que una visión distribucional sea más informativa que una estimación escalar de la media para comprender el comportamiento del agente en torno a la frontera de decisión PERMIT/BLOCK.

\subsection{Factor de descuento $\gamma = 0$: formulación de un solo paso}
En todas las ejecuciones de esta tesis, incluida la ejecución oficial, el factor de descuento se fija en $\gamma = 0$ (valor registrado en el artefacto de configuración de la ejecución). La elección es deliberada y coherente con la formulación conjunto-de-datos-como-entorno: dado que el entorno es un conjunto de datos estático y las acciones del agente no afectan a los estados futuros, no existe dependencia temporal entre flujos que un factor de descuento positivo pudiera explotar.

Con $\gamma = 0$, el error de diferencia temporal definido más arriba se reduce a
\[
    \delta_{ij} = r - \theta_i(s, a),
\]
es decir, el término de bootstrapping $\gamma\,\theta_j^-(s', a^*)$ se anula y la red objetivo deja de contribuir. Cada cuantil aprendido $\theta_i(s,a)$ regresa por tanto hacia la recompensa inmediata de la decisión, y QRDQN actúa como un estimador distribucional de un solo paso: formalmente, un \emph{bandit contextual sensible al coste}, y no un proceso de decisión de Markov secuencial.

Esta reducción no vacía de contenido la elección de QRDQN. La cabeza distribucional sigue modelando la distribución completa de la recompensa asimétrica en torno a la frontera \texttt{PERMIT}/\texttt{BLOCK} —y no únicamente su media—, lo que aporta información sobre la incertidumbre del valor bajo coste asimétrico. La diferencia frente a un clasificador supervisado no reside en la presencia de dinámica secuencial, inexistente aquí por diseño, sino en que el aprendizaje se realiza mediante valores sobre una matriz de coste explícita (con la penalización de falso negativo muy superior a la de falso positivo) en lugar de mediante entropía cruzada sobre etiquetas. La tesis no reivindica autonomía secuencial ni respuesta adversarial en múltiples pasos: tal y como se recoge en las limitaciones de la formulación, esas capacidades quedan fuera del alcance implementado.

\subsection{Estrategia de exploración}
Durante el entrenamiento, el agente emplea una política $\varepsilon$-greedy. Con probabilidad $\varepsilon$ selecciona una acción uniformemente aleatoria de $\{\texttt{PERMIT}, \texttt{BLOCK}\}$; con probabilidad $1-\varepsilon$ selecciona la acción cuya estimación media de cuantiles es más alta \parencite{Mnih2015DQN}. $\varepsilon$ decae linealmente desde un valor inicial cercano a uno hasta un valor final pequeño a lo largo de una fracción fija del total de pasos de entrenamiento, y luego permanece constante. En el entorno de acción binaria, el presupuesto de exploración es relativamente barato porque solo existen dos alternativas. La fracción de exploración y el valor final de $\varepsilon$ se almacenan en el artefacto de configuración de la ejecución para que las ejecuciones con diferentes programas no se mezclen o comparen en silencio sin tener en cuenta sus diferencias de exploración.

\subsection{Configuración del entrenamiento}
La implementación utiliza el módulo QRDQN de \texttt{sb3-contrib} \parencite{SB3ContribQRDQNDocs}. El modelo emplea una política de perceptrón multicapa sobre las observaciones de 152 dimensiones. El entrenamiento utiliza un buffer de repetición, actualizaciones por minibatch, una red objetivo y predicción determinista durante la evaluación. El script de entrenamiento mantenido admite preajustes rápidos y completos, modos de partición aleatoria y por día, semillas explícitas, límites de filas configurables y pasos temporales configurables.

Los valores por defecto del preajuste completo actual utilizan lotes más grandes y más pasos de gradiente que el preajuste rápido. Esta separación permite al proyecto ejecutar comprobaciones rápidas de humo sin confundirlas con las ejecuciones finales del benchmark. Cualquier resultado reportado debe incluir el identificador exacto de la ejecución y la configuración.

\subsection{Procedencia de los hiperparámetros}
Los hiperparámetros de la ejecución oficial ($\gamma = 0$, arquitectura \texttt{[1024, 1024, 512]}, $200$ cuantiles, tasa de aprendizaje $5\times10^{-5}$, $20$ pasos de gradiente y $3\,000\,000$ de pasos de entrenamiento) constituyen un \emph{perfil fijo establecido manualmente}, no el resultado de una búsqueda automática. El repositorio incluye una sonda de optimización con Optuna (\texttt{src/tune\_hparams.py}) con fines exploratorios, pero su espacio de búsqueda —$\gamma \in [0{,}95,\,0{,}999]$, arquitecturas en $\{[256,128],\,[512,256],\,[256,256]\}$ y pasos de gradiente en $\{10,50,100\}$— \emph{no contiene} la configuración oficial. Por tanto, dicha configuración no pudo originarse en esa búsqueda y no debe interpretarse como un óptimo ajustado por Optuna; responde a un diseño manual coherente con la formulación de un solo paso, y queda congelada y verificada por las pruebas del proyecto.

\subsection{Persistencia del modelo y el preprocesamiento}
El entrenamiento produce un directorio de ejecución y un artefacto de modelo. El directorio de ejecución almacena la configuración, las métricas, el escalador ajustado y los límites de percentiles de entrenamiento. Persistir los artefactos de preprocesamiento es esencial porque la inferencia en la Fase 2 debe utilizar el mismo estado de normalización que el entrenamiento. Recalcular un escalador sobre el tráfico de laboratorio filtraría información de la distribución de inferencia hacia la interfaz del modelo y haría la ejecución irreproducible.

\section{Marco de implementación}
El pipeline experimental está implementado en Python y se apoya en un conjunto de bibliotecas de código abierto bien establecidas. El componente de RL profundo principal utiliza la extensión \texttt{sb3-contrib} de Stable Baselines3, que proporciona una implementación contrastada de QRDQN con arquitectura de red configurable, buffer de repetición, programa de exploración y actualizaciones de la red objetivo \parencite{SB3ContribQRDQNDocs}. El envoltorio del entorno sigue la API de Gymnasium \parencite{GymnasiumDocs}, haciendo que el conjunto de datos de flujos sea compatible con cualquier bucle de entrenamiento que cumpla con Gymnasium sin cambios en la implementación del algoritmo.

\subsection{Pila de procesamiento de datos y aprendizaje supervisado}
La carga y el preprocesamiento de datos se apoyan en \texttt{pandas} para la ingesta de CSV, la estandarización de columnas y la lectura fragmentada de ficheros grandes. Las matrices de características se convierten a arrays de \texttt{numpy} y se castean a \texttt{float32} antes de entrar en el entorno de Gymnasium o en la línea base supervisada. La estandarización utiliza \texttt{StandardScaler} de \texttt{scikit-learn}, elegido por su compatibilidad con el protocolo de serialización de \texttt{sklearn}, que permite persistir y recargar los escaladores ajustados sin reimplementación. La línea base de Random Forest también utiliza \texttt{scikit-learn}, consumiendo los mismos arrays de entrada que entran en el entorno de RL. Esta ruta de datos compartida garantiza que cualquier diferencia de rendimiento entre los dos modelos refleje el comportamiento algorítmico y no una discrepancia en el manejo de los datos.

\subsection{Gestión de artefactos}
Cada ejecución de entrenamiento produce un \texttt{RUN\_ID} único compuesto por una marca temporal y un sufijo aleatorio corto. El directorio de ejecución reúne la política QRDQN serializada, el escalador ajustado, los límites de percentiles, un JSON de configuración y un JSON de métricas. Las ejecuciones de validación y de inferencia de la Fase 2 producen subdirectorios independientes con sus propios registros de configuración y métricas. Este esquema de directorios plano evita dependencias de bases de datos al tiempo que hace las ejecuciones recuperables de forma independiente y comparables mediante inspección del sistema de ficheros.

\subsection{Controles de reproducibilidad}
Las semillas se fijan para los generadores de números aleatorios de \texttt{numpy}, \texttt{random} y \texttt{torch} al comienzo de cada ejecución. Los reinicios del entorno de Gymnasium también reciben una semilla explícita. El pipeline está preparado para que, cuando se utilicen múltiples semillas para la estimación de la varianza, cada semilla genere un directorio de ejecución independiente de modo que los resultados por semilla se preserven individualmente en lugar de promediarse en silencio; no obstante, los resultados reportados en esta tesis corresponden a una única semilla y no incluyen una estimación de varianza entre ejecuciones. Estos controles buscan reducir la varianza interna a la ejecución lo suficiente como para distinguir diferencias de rendimiento significativas del ruido, aunque no pueden eliminar todo el indeterminismo derivado del ordenamiento de los kernels de la GPU.
